% id: 130-06
% book: 베이직쎈 공통수학1 · 08 일차부등식
% section: 실전 감각 UP
% figure: tikz:fig-130-06
% answer: ⑤
% answer_source: 답지
% variant_level: 0 (원본 전사)
% 이 파일은 build.mjs 가 items.json 에서 만든다. 고칠 때는 items.json 을 고친다.
\smash{\raisebox{1.5mm}{\dmkichul{베이직쎈 공통수학1 130쪽 06번}}}
\noindent\begin{minipage}[t]{0.58\linewidth}\vspace{0pt}\raggedright 오른쪽 표는 두 식품 $\mathrm{A}$, $\mathrm{B}$의 $1$개당 열량과 단백질의 양을 표로 나타낸 것이다. $\mathrm{A}$, $\mathrm{B}$를 합하여 $20$개의 식품을 섭취하고 $1500$~kcal 이하의 열량과 $45$~g 이하의 단백질을 얻으려고 한다. 이때 식품 $\mathrm{A}$의 개수가 될 수 \underline{없는} 것은?\end{minipage}\hfill\begin{minipage}[t]{0.4\linewidth}\vspace{0pt}\centering \resizebox{\ifdim\width>\linewidth\linewidth\else\width\fi}{!}{% 실전 감각 UP 130-06(본책 130쪽) — 두 식품 A, B 의 1개당 열량(kcal)·단백질(g) 표. 스캔 표를 tabular 로 다시 조판(원문 수치만).
\begin{tabular}{|c|c|c|}
\hline
 & \begin{tabular}[c]{@{}c@{}}열량\\(kcal)\end{tabular} & \begin{tabular}[c]{@{}c@{}}단백질\\(g)\end{tabular} \\
\hline
$\mathrm{A}$ & $50$ & $2.5$ \\
\hline
$\mathrm{B}$ & $100$ & $1.5$ \\
\hline
\end{tabular}
}\end{minipage}\par
\begin{choices32}{$12$}{$13$}{$14$}{$15$}{$16$}\end{choices32}
